\section{Introduction}
\label{sec:introduction}

Competition among suppliers is a standard prescription in public procurement. A larger set of bidders tends to lower the minimum production cost, reduce procurement payments, and improve productive efficiency conditional on a project being undertaken. Those gains are central to auction theory and to procurement policy. Yet procurement competition also changes the incidence of surplus. In particular, stronger competition can compress the information rent earned by the winning supplier. This paper asks whether that incidence effect can feed back into a different policy margin: whether a regional government is willing to undertake the supplier-side activity required for an interregional industrial project.

The question matters whenever industrial policy is organized around complementary roles rather than a single undifferentiated subsidy. Governments increasingly use procurement, local-content initiatives, supply-chain policies, and targeted industrial support to build or localize complementary stages of production. Public procurement is itself recognized as an industrial-policy instrument, including in innovation policy \citep{ChiappinelliGiuffridaSpagnolo2025}. The OECD's recent review documents procurement measures used to support domestic and local production and emphasizes the tension among industrial-policy objectives, market access, competition, and supply-chain resilience \citep{OECD2026}. Territorial procurement is also empirically relevant: local and devolved governments differ systematically in the extent to which their procurement is sourced within their own regions \citep{EckersleyEtAl2023}. Recent industrial-organization work on industrial policy likewise emphasizes that policy effects depend on market structure, that localizing supply chains can be central, and that government policy may have a coordination role across market actors \citep{VanBiesebroeckZhang2026,MottaPolo2026}. In such settings, a policy environment with more supplier competition can improve production efficiency while changing the local rents that make a jurisdiction willing to host the supplier-side activity.

We study the smallest model that isolates this interaction. Two regional governments simultaneously choose between a locally attractive activity, denoted $U$, and a complementary supplier-side activity, denoted $D$. The direct local return to $U$ exceeds that to $D$ by $\Delta>0$. If the regions choose complementary roles, $(U,D)$ or $(D,U)$, the $U$ region is the lead or buyer side of the complementary project and the $D$ region hosts or develops the supplier-side industrial base. The $D$ side then hosts a procurement market with $N\geq2$ symmetric suppliers. Suppliers have independent private costs and compete for a project whose value is $v$. The baseline uses a second-price reverse auction because it makes the surplus decomposition transparent: expected total match surplus depends on the lowest cost, buyer surplus depends on the second-lowest cost, and winning supplier rent is the gap between the second- and first-order statistics.

The lower-stage comparative static is deliberately standard. Let $C_{1:N}$ and $C_{2:N}$ denote the lowest and second-lowest supplier costs. Define
\begin{equation}
G_N\equiv v-\E[C_{1:N}],
\qquad
R_N\equiv \E[C_{2:N}-C_{1:N}].
\label{eq:intro-GR}
\end{equation}
The first object is expected real surplus from an active complementary project and the second is expected winning supplier information rent. More suppliers reduce the expected minimum cost, so $G_N$ rises. Under the standard decreasing-reversed-hazard-rate (DRHR) condition used in procurement theory, expected winning rent $R_N$ strictly falls with the number of suppliers and converges to zero \citep{Li2005,XuLi2019}; related bidder-attraction results appear in \citet{Szech2011}. We do not claim this rent-dissipation result as new.

The new result begins when these two effects enter the regional policy game. Suppose the $D$ region internalizes a fraction $\rho\in(0,1]$ of local supplier rent. A coordinated policymaker values the full real surplus $G_N$ created by regional differentiation. The $D$ government, however, gives up its direct premium $\Delta$ and obtains only $\rho R_N$ from the supplier-side role. Let $A>0$ scale the mass or importance of the complementary project. The coordinated differentiation threshold is
\begin{equation}
A_N^P=\frac{\Delta}{G_N},
\label{eq:intro-AP}
\end{equation}
whereas the decentralized threshold at which a region is willing to accept $D$ against a region choosing $U$ is
\begin{equation}
A_N^N=\frac{\Delta}{\rho R_N}.
\label{eq:intro-AN}
\end{equation}
At a fixed $N$, this is an incidence wedge: the policymaker who values the full project switches earlier than the jurisdiction that captures only part of its benefit. That fixed-market observation is not the contribution of the present paper. It maps directly into the familiar logic that incomplete local capture delays decentralized adjustment.

Our contribution is the comparative static in market structure. Because supplier competition raises $G_N$ but lowers $R_N$, it moves the two thresholds in opposite directions:
\begin{equation}
A_{N+1}^P<A_N^P,
\qquad
A_{N+1}^N>A_N^N.
\label{eq:intro-divergence}
\end{equation}
Competition therefore makes differentiation socially desirable at a lower project scale while making the supplier-side region require a larger project before voluntarily accepting that role. The set of project scales for which decentralized policy uniquely duplicates $U$ even though coordinated policy differentiates consequently expands with supplier competition.

The stronger result concerns the equilibrium set rather than a threshold gap. For a fixed $N$, the induced government game is an anti-coordination game only when $A\rho R_N>\Delta$. In that region, $(U,D)$ and $(D,U)$ are differentiated pure Nash equilibria. Since $R_N\to0$ as supplier competition becomes sufficiently strong, those differentiated equilibria eventually disappear. At the same time, the planner's valuation of differentiation rises because $G_N$ increases. If
\begin{equation}
A(v-c_0)>\Delta,
\label{eq:intro-limit-condition}
\end{equation}
where $c_0$ is the lower endpoint of the cost support, then sufficiently large $N$ yields
\begin{equation}
\text{coordinated optimum}=\{(U,D),(D,U)\},
\qquad
\text{unique Nash equilibrium}=(U,U).
\label{eq:intro-reversal}
\end{equation}
Thus competition can destroy efficient differentiated equilibria and leave priority duplication as the unique decentralized outcome. We call this \emph{competition-induced policy reversal}.

The result is counterintuitive because the supplier market becomes more efficient conditional on the project being undertaken. The failure arises at the policy-selection layer. More competition reduces production cost but also reallocates surplus away from the local supplier side. If local governments respond to locally captured income or rents, the same increase in supplier-market competition that improves procurement performance can weaken the incentive of one region to specialize in the activity required for interregional complementarity. The policy problem is therefore not that competition directly lowers welfare. It is that competition changes the incidence supporting decentralized coordination.

The paper connects four literatures. First, auction and procurement theory studies bidder competition, participation, rent, and procurement preferences. \citet{Li2005} establishes monotonicity of expected rent under DRHR; \citet{Szech2011} studies costly bidder attraction and the difference between private and social bidder-count objectives; and \citet{XuLi2019} applies related order-statistic results to procurement competition and welfare. Endogenous participation can overturn simple competition intuitions, as \citet{LiZheng2009} show, which is why the number of suppliers is a market-structure primitive in our baseline rather than an endogenous entry choice. Procurement preferences and set-asides also redistribute surplus and alter participation \citep{Vagstad1995,Marion2007,AtheyCoeyLevin2013}. Regional procurement restrictions can explicitly trade off participation against local-industry objectives \citep{SaitoMitsuta2012}, local-content policy design can depend on local-supplier competitiveness \citep{Macatangay2016}, and actual local procurement displays a territorial sourcing dimension \citep{EckersleyEtAl2023}. Our paper takes these procurement-side forces as inputs and asks how they change a strategic public-policy game.

Second, industrial-organization research increasingly studies industrial policy as a function of market structure. \citet{AghionEtAl2015} document an interaction between industrial policy and competition, while the recent IJIO industrial-policy agenda emphasizes market structure, supply-chain localization, and coordination \citep{VanBiesebroeckZhang2026}. \citet{MottaPolo2026} provide a recent theoretical example in which market organization and appropriability shape precautionary industrial policy under supply-chain disruption. Our mechanism is different: supplier competition changes local rent incidence and thereby alters which region is willing to take a complementary policy role.

Third, place-based and local-government policy can itself be understood through competition, differentiation, and rent allocation. In an earlier regional-competition model, \citet{JustmanThisseVanYpersele2005} show that regions may differentiate infrastructure to soften competition for investors and that informational frictions can make fiscal competition inefficient. \citet{Slattery2025} models state and local governments bidding for firms and shows that intergovernmental competition can transfer rents to firms. \citet{LinLi2026} study local industrial-policy competition in a quantitative spatial equilibrium environment. We study a different direction of strategic influence: competition among suppliers changes the equilibrium of a game among governments, and can eliminate differentiated regional-policy equilibria rather than generate differentiation as a direct strategic response between regions.

Finally, the paper is related to a companion analysis of fixed-capacity regional industrial-policy composition \citep{Companion2026}. That paper takes complementary surplus and its local incidence largely as primitives and studies a continuous portfolio game. Here the policy choice is binary and the incidence wedge is generated by an explicit private-information procurement market. At any fixed number of suppliers, the unilateral switching condition can be mapped to the earlier reduced-form logic; we do not treat that fixed-market comparison as a new result. The new object is the supplier-market comparative static: stronger competition changes both total surplus and local rent capture and can alter the exact Nash-equilibrium set of regional policy.

The remainder of the paper proceeds as follows. Section \ref{sec:institutional} motivates the procurement--industrial-policy link. Section \ref{sec:model} presents the model and procurement continuation. Section \ref{sec:equilibrium} characterizes decentralized regional-policy equilibrium, and Section \ref{sec:planner} defines the coordinated fixed-instrument benchmark. Section \ref{sec:main-results} establishes competition--coordination divergence and competition-induced policy reversal. Section \ref{sec:robustness} studies incidence, trade failure, and auction-format robustness. Section \ref{sec:discussion} develops empirical implications and scope conditions. Section \ref{sec:conclusion} concludes. Proofs are in the Appendix.
